% !TeX root = ../../Book1.tex
\section{Doob 最大不等式}
\label{b1:sec:3104}

期望和范数通常控制一个确定时刻的大小。
若关心过程在观察期间是否曾经越过某个阈值，就必须控制整段路径的最大值。
把每个时刻的 Markov 不等式相加会得到随观察次数增长的上界；
鞅的条件平均关系能够把各个越界时刻重新联系到同一个终点。
本节的关键是按第一次越界的时刻拆分事件，随后仍只在确定时刻之间使用塔性质。

\subsection{有限时段的最大过程}
\label{b1:subsec:310401}

设 \((\Omega,\mathcal F,\mathbb P)\) 为概率空间，
\((\mathcal F_n)_{0\leq n\leq N}\) 为其上的滤过，
\((X_n)_{0\leq n\leq N}\) 为关于该滤过的非负实值适应过程。
它在时刻 \(n\) 以前的\term[zuida-guocheng]{最大过程}[maximal process]记为
\[
 X_n^*\coloneqq\max_{0\leq k\leq n}X_k.
\]
对于带符号的实过程 \((M_n)\)，通常关心的是
\(\max_{0\leq k\leq n}|M_k|\)，这就是非负过程 \((|M_k|)\) 的最大过程。
有限最大值对 \(\mathcal F_n\) 可测，因为每个 \(X_k\)、\(k\leq n\)
都对 \(\mathcal F_n\) 可测。若各项可积且非负，则
\(X_n^*\leq\sum_{k=0}^nX_k\)，所以有限时段的最大值也可积。
\symidx{\(X_n^*\)：非负过程在时段 \(0,\ldots,n\) 上的最大值}

\ProseParagraphBreak

最大值可能出现得很早，终点值却已经回落。
因此没有逐点不等式 \(X_N^*\leq C X_N\) 可供直接积分。
例如一个非负公平资本过程可能先翻倍，随后输掉全部资本。
最大不等式比较的是这些变量的概率或范数，而非每一条路径上的终点和最高点。

\ProseParagraphBreak

为了看出所需的组织方式，固定 \(\lambda>0\)，写出
\[
 A_k=\{X_0<\lambda,\ldots,X_{k-1}<\lambda,\ X_k\geq\lambda\},
 \quad 0\leq k\leq N.
\]
当 \(k=0\) 时，前面的条件为空，故 \(A_0=\{X_0\geq\lambda\}\)。
这些事件两两不交，且
\[
 \{X_N^*\geq\lambda\}=\bigcup_{k=0}^N A_k.
\]
每个 \(A_k\) 属于 \(\mathcal F_k\)：它只问到时刻 \(k\) 为止的取值，
而其中每个 \(X_j\)、\(j\leq k\) 都对 \(\mathcal F_k\) 可测。
这个可测性使我们可以在 \(A_k\) 上比较 \(X_k\) 与 \(X_N\) 的积分。
若按“全程最大值在哪个时刻取得”拆分，事件却可能依赖未来取值，
便不能直接作同样的条件期望计算。

\subsection{弱型最大估计}
\label{b1:subsec:310402}

下面的\term[doob-zuida-budengshi]{Doob 最大不等式}[Doob maximal inequality]%
先给出一次矩控制下的越界概率。
保留终值在越界事件上的积分，比直接用 \(\mathbb E X_N\) 代替它更精确；
后面强型估计需要的正是这项局部化信息。

\begin{theorem}[Doob 弱型最大不等式]
\label{b1:thm:doob-weak-31}
设 \(N\geq0\) 为整数，\((\Omega,\mathcal F,\mathbb P)\) 为概率空间，
\((\mathcal F_n)_{0\leq n\leq N}\) 为滤过，
\((X_n)_{0\leq n\leq N}\) 为关于该滤过的非负实值下鞅。
令 \(X_N^*=\max_{0\leq k\leq N}X_k\)。则对每个 \(\lambda>0\)，有
\begin{equation}
 \lambda\cdot\mathbb P(X_N^*\geq\lambda)
 \leq\mathbb E\bigl(X_N\cdot\mathbf1_{\{X_N^*\geq\lambda\}}\bigr)
 \leq\mathbb E X_N.
 \label{b1:eq:doob-local-tail-31}
\end{equation}
将两个事件中的 \(\geq\lambda\) 同时换成 \(>\lambda\)，结论也成立。
\end{theorem}
\begin{proof}
按上一小节的公式定义首次达到阈值的事件 \(A_k\)。
在 \(A_k\) 上有 \(X_k\geq\lambda\)，所以
\[
 \lambda\cdot\mathbb P(A_k)
 \leq\mathbb E\bigl(X_k\cdot\mathbf1_{A_k}\bigr).
\]
由\cref{b1:prop:martingale-multistep-31}，
\(X_k\leq\mathbb E(X_N\mid\mathcal F_k)\)。
又因 \(A_k\in\mathcal F_k\)，条件期望的定义给出
\[
 \begin{aligned}
 \mathbb E\bigl(X_k\cdot\mathbf1_{A_k}\bigr)
 &\leq\mathbb E\bigl(\mathbb E(X_N\mid\mathcal F_k)
                         \cdot\mathbf1_{A_k}\bigr)\\
 &=\mathbb E\bigl(X_N\cdot\mathbf1_{A_k}\bigr).
 \end{aligned}
\]
对 \(k=0,\ldots,N\) 求和。事件两两不交，且其并为
\(\{X_N^*\geq\lambda\}\)，故得第一个不等式。
第二个不等式用到 \(X_N\geq0\)。
若把首次越界事件改成
\(\{X_0\leq\lambda,\ldots,X_{k-1}\leq\lambda,X_k>\lambda\}\)，
完全相同的计算给出严格越界版本。
\end{proof}

证明没有把随机选择的时刻直接代入下鞅定义。
它对每个确定的 \(k\) 分别作条件平均，然后利用可测事件分块相加。
这一点使首次越界方法只依赖前面已经建立的条件期望性质。

\begin{corollary}[鞅的最大概率估计]
\label{b1:cor:martingale-maximal-31}
设 \(N\geq0\) 为整数，\((M_n)_{0\leq n\leq N}\) 为概率空间 \((\Omega,\mathcal F,\mathbb P)\)
上关于滤过 \((\mathcal F_n)\) 的实值或复值鞅。
则对每个 \(\lambda>0\)，有
\[
 \mathbb P\left(\max_{0\leq k\leq N}|M_k|\geq\lambda\right)
 \leq\frac{\mathbb E|M_N|}{\lambda}.
\]
若另有 \(M_N\in L^r(\mathbb P)\)、\(1<r<\infty\)，则
\[
 \mathbb P\left(\max_{0\leq k\leq N}|M_k|\geq\lambda\right)
 \leq\frac{\mathbb E|M_N|^r}{\lambda^r}.
\]
\end{corollary}
\begin{proof}
由\cref{b1:prop:martingale-convex-31}，\((|M_n|)\) 是非负下鞅，
故第一式由弱型最大不等式得到。
对于第二式，先用 \(M_k=\mathbb E(M_N\mid\mathcal F_k)\)
和条件期望的 \(L^r\) 收缩性，确认每个 \(M_k\in L^r\)。
非负下鞅 \((|M_k|)\) 经非减凸函数 \(t\mapsto(t^+)^r\) 作用后，
仍是非负下鞅 \((|M_k|^r)\)。把弱型最大不等式用于此过程及阈值 \(\lambda^r\)，
得到第二式。
\end{proof}

单独在终点应用 Markov 不等式只控制 \(\{|M_N|\geq\lambda\}\)，
上述推论同时控制先前任意时刻发生的越界。
例如，若 \(Z\in L^1(\mathbb P)\)，并对任意滤过取
\(M_n=\mathbb E(Z\mid\mathcal F_n)\)，则
\[
 \mathbb P\left(\max_{0\leq k\leq N}
       |\mathbb E(Z\mid\mathcal F_k)|\geq\lambda\right)
 \leq\frac{\mathbb E|Z|}{\lambda}.
\]
右端不随信息更新的次数 \(N\) 增长。

\subsection{强型最大估计}
\label{b1:subsec:310403}

弱型估计说明阈值越高，越界概率越小；强型估计则把所有阈值的贡献合在一起，
控制最大过程的 \(L^p\) 范数。
直接积分粗略上界 \(\mathbb P(X_N^*\geq\lambda)\leq\mathbb E X_N/\lambda\)
不会得到所需结果。必须先保留\eqref{b1:eq:doob-local-tail-31}中
随阈值变化的事件，再交换非负积分。

\begin{theorem}[Doob 强型最大不等式]
\label{b1:thm:doob-strong-31}
设 \(N\geq0\) 为整数，\((\Omega,\mathcal F,\mathbb P)\) 为概率空间，
\((\mathcal F_n)_{0\leq n\leq N}\) 为滤过，
\((X_n)_{0\leq n\leq N}\) 为关于该滤过的非负实值下鞅。
若 \(1<p<\infty\) 且 \(X_N\in L^p(\mathbb P)\)，则
\[
 \left\|\max_{0\leq k\leq N}X_k\right\|_p
 \leq\frac p{p-1}\cdot\|X_N\|_p.
\]
若 \(X_N\in L^\infty(\mathbb P)\)，则
\[
 \left\|\max_{0\leq k\leq N}X_k\right\|_\infty
 \leq\|X_N\|_\infty.
\]
特别地，若 \((M_n)_{0\leq n\leq N}\) 是关于该滤过的实值或复值鞅，
\(1<p<\infty\) 且 \(M_N\in L^p(\mathbb P)\)，则
\[
 \left\|\max_{0\leq k\leq N}|M_k|\right\|_p
 \leq\frac p{p-1}\cdot\|M_N\|_p
 \quad(1<p<\infty),
\]
若 \(M_N\in L^\infty(\mathbb P)\)，则
\(\|\max_{k\leq N}|M_k|\|_\infty\leq\|M_N\|_\infty\)。
\end{theorem}
\begin{proof}
记 \(Y=X_N^*\)。先确认所需的范数有限，以免在最后消去一个无穷量。
对每个 \(k\leq N\)，下鞅的多步比较给出
\[
 0\leq X_k\leq\mathbb E(X_N\mid\mathcal F_k).
\]
条件期望的 \(L^p\) 收缩性说明 \(X_k\in L^p\)，
而 \(Y^p\leq\sum_{k=0}^N X_k^p\) 说明 \(Y\in L^p\)。

对任意非负有限实数 \(y\)，有
\(y^p=\int_0^y p\lambda^{p-1}\,\mathrm d\lambda\)。
在 \(\Omega\times(0,\infty)\) 上，\(\mathbb P\) 与 Lebesgue 测度都为
\(\sigma\)-有限测度，且下列被积函数非负并联合可测。
用\cref{b1:thm:tonelli}交换积分，再用局部尾部估计，得到
\[
 \begin{aligned}
 \mathbb E Y^p
 &=\int_0^\infty p\lambda^{p-1}\cdot
                    \mathbb P(Y\geq\lambda)\,\mathrm d\lambda\\
 &\leq p\int_0^\infty\lambda^{p-2}\cdot
        \mathbb E\bigl(X_N\cdot\mathbf1_{\{Y\geq\lambda\}}\bigr)
                          \,\mathrm d\lambda\\
 &=p\cdot\mathbb E\left(X_N\int_0^Y\lambda^{p-2}\,\mathrm d\lambda\right)
 =\frac p{p-1}\cdot\mathbb E(X_NY^{p-1}).
 \end{aligned}
\]
在零点附近积分 \(\lambda^{p-2}\) 有限，正好需要 \(p>1\)。
令 \(q=p/(p-1)\)。Hölder 不等式给出
\[
 \mathbb E(X_NY^{p-1})
 \leq\|X_N\|_p\cdot\|Y^{p-1}\|_q
 =\|X_N\|_p\cdot\|Y\|_p^{p-1}.
\]
若 \(\|Y\|_p=0\)，结论已经成立；否则两边除以有限正数
\(\|Y\|_p^{p-1}\)，即得所需估计。

若 \(X_N\in L^\infty\)，仍用
\(0\leq X_k\leq\mathbb E(X_N\mid\mathcal F_k)\leq\|X_N\|_\infty\)。
有限多个时刻的异常零测集合并后，最大值也几乎必然不超过该常数。
最后将已经证明的结论应用于非负下鞅 \((|M_n|)\)。
\end{proof}

强型估计的常数与时段长度无关，因此可以让 \(N\) 趋于无穷。
这一过程只需要最大值随 \(N\) 单调增加，尚不需要鞅本身已经收敛。

\begin{corollary}[无穷时段的最大估计]
\label{b1:cor:doob-infinite-horizon-31}
设 \((M_n)_{n\geq0}\) 是概率空间 \((\Omega,\mathcal F,\mathbb P)\)
上关于滤过 \((\mathcal F_n)\) 的实值或复值鞅，令
\(M^*=\sup_{n\geq0}|M_n|\)。
若 \(C_1=\sup_n\mathbb E|M_n|<\infty\)，则
\[
 \mathbb P(M^*>\lambda)\leq\frac{C_1}{\lambda}
 \quad(\lambda>0),
\]
并且 \(M^*<\infty\) 几乎必然。
若 \(1<p<\infty\) 且 \(C_p=\sup_n\|M_n\|_p<\infty\)，则
\[
 \|M^*\|_p\leq\frac p{p-1}\,C_p.
\]
\end{corollary}
\begin{proof}
事件 \(\{\max_{k\leq N}|M_k|>\lambda\}\) 随 \(N\) 增加，
其并为 \(\{M^*>\lambda\}\)。由测度的下连续性和有限时段弱型估计，得到第一式。
再令整数阈值趋于无穷，测度的上连续性给出
\(\mathbb P(M^*=\infty)=0\)。
对强型估计，\((\max_{k\leq N}|M_k|)^p\) 单调增加到 \((M^*)^p\)，
故单调收敛定理和有限时段估计给出
\[
 \mathbb E(M^*)^p
 =\lim_{N\to\infty}\mathbb E\left(\max_{k\leq N}|M_k|\right)^p
 \leq\left(\frac p{p-1}\right)^p C_p^p.
\]
\end{proof}

一次矩一致有界的鞅因此几乎每条路径都有有限上界。
有界数列仍可能反复振荡，所以上述结论还不能代替收敛定理。
下一节将专门估计路径穿过固定区间的次数。

\subsection{最大估计的应用与边界}
\label{b1:subsec:310404}

先看二阶矩估计的具体意义。
设 \(\xi_1,\ldots,\xi_N\) 为独立、平方可积、均值为零、方差分别为
\(\sigma_1^2,\ldots,\sigma_N^2\) 的实随机变量，令
\(S_n=\sum_{j=1}^n\xi_j\)、\(S_0=0\)。
独立和构成鞅，而\cref{b1:prop:30-variance-sum}给出
\(\mathbb E|S_N|^2=\sum_{j=1}^N\sigma_j^2\)。于是
\[
 \begin{aligned}
 \mathbb E\left(\max_{0\leq k\leq N}|S_k|\right)^2
 &\leq4\sum_{j=1}^N\sigma_j^2,\\
 \mathbb P\left(\max_{0\leq k\leq N}|S_k|\geq\lambda\right)
 &\leq\frac1{\lambda^2}\sum_{j=1}^N\sigma_j^2.
 \end{aligned}
\]
第一项来自强型估计，第二项直接把弱型估计用于 \((S_k^2)\)。
若先对第一项应用 Markov 不等式，概率上界会多出常数 \(4\)；
这个比较提醒我们，控制同一对象时也应选择合适的版本。

\ProseParagraphBreak

最大估计还把终值的逼近转化为整列条件平均的同步逼近。
设 \(Z,W\in L^p(\mathbb P)\)、\(1<p<\infty\)，
对任意固定滤过取 \(M_n=\mathbb E(Z\mid\mathcal F_n)\)、
\(L_n=\mathbb E(W\mid\mathcal F_n)\)。差过程也是鞅，且
\(\|M_n-L_n\|_p\leq\|Z-W\|_p\)，所以
\[
 \left\|\sup_{n\geq0}|M_n-L_n|\right\|_p
 \leq\frac p{p-1}\cdot\|Z-W\|_p.
\]
当 \(p=1\) 时仍有相应的概率控制
\[
 \mathbb P\left(\sup_{n\geq0}|M_n-L_n|>\lambda\right)
 \leq\frac{\|Z-W\|_1}{\lambda}.
\]
这两式都控制全部时刻的误差，而不只是固定 \(n\) 的误差。

\ProseParagraphBreak

强型估计中的 \(p>1\) 不能改成 \(p\geq1\)。
下面给出完全有限的计算，说明不存在与时段长度无关的常数 \(C\)，
使所有非负鞅都满足 \(\mathbb E M_N^*\leq C\mathbb E M_N\)。
取独立公平符号列 \((\xi_k)\)，令
\[
 M_0=1,\quad
 M_n=2^n\cdot\mathbf1_{\{\xi_1=\cdots=\xi_n=1\}},
 \quad\mathcal F_n=\sigma(\xi_1,\ldots,\xi_n).
\]
在连续得到 \(1\) 时资本每步翻倍，第一次得到 \(-1\) 后则永远为零。
每个 \(M_n\) 有界，因此可积；由独立性，
\[
 \mathbb E(M_{n+1}\mid\mathcal F_n)
 =2^{n+1}\cdot\mathbf1_{\{\xi_1=\cdots=\xi_n=1\}}\cdot\frac12=M_n.
\]
所以这是非负鞅，且 \(\mathbb E M_N=1\)。
在前 \(N\) 次中，恰有开头 \(r\) 次为 \(1\)、随后一次为 \(-1\)
的事件概率为 \(2^{-(r+1)}\)，其中 \(0\leq r<N\)；该事件上的最大值为 \(2^r\)。
若前 \(N\) 次全为 \(1\)，概率为 \(2^{-N}\)，最大值为 \(2^N\)。
因此
\begin{equation}
 \mathbb E\left(\max_{0\leq k\leq N}M_k\right)
 =\sum_{r=0}^{N-1}2^r\cdot2^{-(r+1)}+2^N\cdot2^{-N}
 =1+\frac N2.
 \label{b1:eq:doob-endpoint-counterexample-31}
\end{equation}
终点的一次矩恒为 \(1\)，全程最大值的一次矩却随 \(N\) 无界增长。
这与弱型估计并不冲突：对整数 \(1\leq j\leq N\)，还有精确等式
\[
 \mathbb P\left(\max_{0\leq k\leq N}M_k\geq2^j\right)=2^{-j}.
\]
越界概率具有一次幂衰减，却不能给出统一的一次矩最大上界。%
\secref{b1:sec:3106}会从一致可积性的角度继续解释同一类资本集中现象。
